\chapter{Use Case 1: Vergleich von Schwachstellen-Scannern}
\label{chap:uc1}

Dieses Kapitel setzt den ersten Use Case der Forschungsumgebung um und wertet ihn aus. Es ist als
geschlossene Einheit aufgebaut und führt von der Zielsetzung über die Konzeption und die
Evaluationsmethodik bis zur Implementierung, den Ergebnissen und deren Diskussion.

\section{Zielsetzung und Bezug zu den Forschungsfragen}
\label{sec:uc1-ziel}

Use Case 1 vergleicht mehrere Schwachstellen-Scanner unter identischen und reproduzierbaren
Bedingungen. Als Ziel dient ein absichtlich verwundbares System, das in der ausgewählten Plattform
CyberRangeCZ (Kapitel~\ref{chap:auswahl}) als isolierte Sandbox bereitgestellt wird. Verglichen
werden ein klassischer quelloffener Scanner (OpenVAS), ein moderner template-basierter Scanner
(Nuclei) und ein kommerzielles Werkzeug (Nessus). Ergänzend kommen Werkzeuge für die Host-Härtung
und die Container-Analyse hinzu.

Der Use Case ist dabei mehr als ein reiner Werkzeugvergleich. Er dient als konkreter Nachweis, dass
die aufgebaute Umgebung reproduzierbare Experimente mit externen Sicherheitswerkzeugen ermöglicht.
Damit operationalisiert er die zentrale Forschungsfrage, die nach einer Forschungsumgebung mit
reproduzierbaren Use Cases fragt.

Inhaltlich ordnet sich der Use Case in den Resilienzbegriff dieser Arbeit ein.
Schwachstellen-Scanning gehört zur vorbereitenden Phase der Resilienz
(Abschnitt~\ref{subsec:resilienz-bezug}), also dem Erkennen von Schwächen, bevor ein Angriff
stattfindet. Je zuverlässiger und vollständiger ein Werkzeug Schwachstellen erkennt, desto besser
ist ein System auf mögliche Angriffe vorbereitet.

Bezogen auf die Untersuchungsfragen aus der Einleitung (Kapitel~\ref{chap:einleitung}) adressiert Use Case 1
vor allem die Modellierung eines reproduzierbaren Use Case (UF3), die Einbindung externer Werkzeuge
über geeignete Schnittstellen (UF4) und die Evaluation der Reproduzierbarkeit (UF5). Bezogen auf die
Bewertungskriterien aus Kapitel~\ref{chap:auswahl} betrifft er vor allem die Reproduzierbarkeit (K1), die
Automatisierung und Lifecycle-Steuerung (K2), die Integrationsfähigkeit externer Werkzeuge (K3)
sowie die Wartbarkeit (K7).

\section{Konzeption}
\label{sec:uc1-konzeption}

Die Konzeption von Use Case 1 umfasst drei Festlegungen: die Auswahl der verglichenen Werkzeuge, die
methodische Einteilung in Untersuchungsperspektiven und das Design des verwundbaren Zielsystems.

\subsection{Auswahl der Werkzeuge}
\label{subsec:uc1-werkzeuge}

Für den Vergleich wurden Werkzeuge ausgewählt, die unterschiedliche Ansätze und Lizenzmodelle
abdecken. Im Zentrum stehen drei Netzwerk-Scanner. OpenVAS steht für den klassischen Ansatz eines
vollständigen Vulnerability-Managements und ist quelloffen. Nuclei ist ein moderner,
leichtgewichtiger und template-basierter Scanner und ebenfalls quelloffen. Nessus ist ein
kommerzielles, proprietäres Werkzeug. Diese Auswahl bildet bewusst ein Spektrum ab, von einer
schweren, umfassenden Plattform über ein schlankes, exploit-nahes Werkzeug bis zu einer
kommerziellen Lösung.

Ein einzelner Netzwerk-Scanner kann jedoch nicht alle relevanten Schwachstellen erfassen. Schwächen
in der Systemkonfiguration oder im Inneren von Containern sind über das Netzwerk nicht sichtbar.
Daher wurden ergänzende Werkzeuge einbezogen. Lynis und OpenSCAP prüfen die Härtung und Compliance
des Systems von innen, Trivy untersucht Container-Images auf bekannte Schwachstellen. Die einzelnen
Werkzeugkategorien wurden bereits in Abschnitt~\ref{sec:vuln-scanner} eingeführt.

\subsection{Untersuchungsperspektiven}
\label{subsec:uc1-perspektiven}

Die ausgewählten Werkzeuge lassen sich nicht sinnvoll auf einer einzigen Skala vergleichen, weil sie
unterschiedliche Fragen beantworten. Statt einer Rangfolge ordnet Use Case 1 die Werkzeuge daher
vier methodisch verschiedenen Untersuchungsperspektiven zu. Ein fairer Vergleich ist nur innerhalb
einer Perspektive möglich. Über die Perspektiven hinweg geht es nicht um eine Rangfolge, sondern um
die Abdeckung unterschiedlicher Sichtweisen auf dasselbe Zielsystem. Tabelle~\ref{tab:uc1-perspektiven}
fasst die vier Perspektiven zusammen.

\begin{table}[H]
\centering
\renewcommand{\arraystretch}{1.3}
\begin{tabular}{p{2.2cm}p{6.0cm}p{4.2cm}}
\hline
\textbf{Perspektive} & \textbf{Frage} & \textbf{Werkzeuge} \\
\hline
A1 -- Unauthenticated Network & Was sieht ein Angreifer von außen? & Nuclei, OpenVAS (unauth) \\
A2 -- Credentialed Network & Was zeigt sich mit Login (Paketstand)? & OpenVAS (credentialed) \\
A3 -- Host Hardening / Compliance & Wie gehärtet ist das System selbst? & Lynis, OpenSCAP/CIS \\
A4 -- Container Image / SCA & Was steckt in den Containern? & Trivy \\
\hline
\end{tabular}
\caption{Die vier Untersuchungsperspektiven von Use Case 1 mit der jeweiligen Leitfrage und den zugeordneten Werkzeugen.}
\label{tab:uc1-perspektiven}
\end{table}

Perspektive A1 betrachtet das System aus der Sicht eines Angreifers ohne Zugangsdaten und umfasst die
beiden Netzwerk-Scanner Nuclei und OpenVAS im unauthentifizierten Betrieb. Perspektive A2 nutzt einen
Login und liest den installierten Paketstand des Hosts aus, was im Use Case über OpenVAS im
credentialed-Modus abgebildet wird. Perspektive A3 bewertet die Härtung und Konfiguration des Systems
von innen und wird durch Lynis und OpenSCAP abgedeckt. Perspektive A4 schaut in die Container-Images
hinein und wird durch Trivy vertreten.

\subsection{Design des Zielsystems}
\label{subsec:uc1-ziel-design}

Damit alle vier Perspektiven etwas zu untersuchen haben, wurde ein absichtlich verwundbares
Zielsystem entworfen. Es wird über Ansible-Rollen reproduzierbar bereitgestellt, sodass sich derselbe
Ausgangszustand wiederholt herstellen lässt. Der genaue technische Aufbau und die Provisionierung
werden in Abschnitt~\ref{sec:uc1-implementierung} beschrieben. Dieser Abschnitt beschreibt die
geplanten Schwachstellen selbst und erläutert, warum sie ausgewählt wurden.

Die Schwachstellen verteilen sich bewusst auf drei Klassen. Die erste Klasse besteht aus
netzwerk-erreichbaren Diensten mit bekannten CVEs. Die zweite Klasse umfasst Fehlkonfigurationen auf
Netzwerkebene. Die dritte Klasse besteht aus host-lokalen Schwächen, die nur mit Zugangsdaten
sichtbar sind. Tabelle~\ref{tab:uc1-zielsystem} gibt einen Überblick, die folgenden Unterabschnitte
beschreiben jede Schwachstelle im Detail.

\begin{table}[H]
\centering
\renewcommand{\arraystretch}{1.25}
\begin{tabular}{p{1.0cm}p{7.5cm}p{3.4cm}}
\hline
\textbf{ID} & \textbf{Dienst / Schwachstelle} & \textbf{Port / Ebene} \\
\hline
A1 & Apache httpd 2.4.49 -- CVE-2021-41773 (Path Traversal) & 8080 \\
A2 & Apache Tomcat 9.0.30 -- CVE-2020-1938 (Ghostcat) & 8009 / 8090 \\
A3 & Grafana 8.3.0 -- CVE-2021-43798 (Path Traversal) & 3000 \\
A4 & Jenkins 2.138.1 -- CVE-2018-1000861 (Pre-Auth-RCE) & 8085 \\
\hline
B1 & FTP mit anonymem Zugang und schwachen Zugangsdaten & 21 \\
B2 & Schwache SSH-Konfiguration & 22 \\
B3 & Anonyme Samba-Freigaben & 139 / 445 \\
B4 & NFS-Export mit \texttt{no\_root\_squash} & 2049 / 111 \\
B5 & Apache-Fehlkonfiguration und Informationsleck & 80 \\
B6 & rlogin-Dienst mit Klartext-Authentisierung & 513 \\
B7 & TFTP-Dienst über UDP & 69/UDP \\
\hline
C1 & World-writable sensible Dateien & Host-lokal \\
C2 & SUID-Bits auf Standardwerkzeugen & Host-lokal \\
C3 & Unsichere sudoers-Einträge (NOPASSWD) & Host-lokal \\
C4 & Schwache lokale Benutzerkonten & Host-lokal \\
\hline
\end{tabular}
\caption{Überblick über die im Zielsystem geplanten Schwachstellen.}
\label{tab:uc1-zielsystem}
\end{table}

\subsubsection{Netzwerk-erreichbare Dienste mit bekannten CVEs}

Die erste Klasse bilden bewusst veraltete Dienste, die jeweils eine dokumentierte CVE-Schwachstelle
mitbringen. Sie sind über das Netzwerk erreichbar und richten sich damit primär an die
unauthentifizierte Perspektive (A1).

Der erste Dienst ist ein Apache-Webserver in der Version 2.4.49 mit CVE-2021-41773. Diese Lücke ist
ein Path-Traversal-Fehler. Über speziell kodierte URLs lassen sich unter bestimmten Bedingungen
Dateien außerhalb des Web-Wurzelverzeichnisses lesen, und bei aktiviertem CGI ist sogar
Codeausführung möglich. Entscheidend ist, dass die Ausnutzbarkeit stark von der Konfiguration
abhängt. Genau deshalb wurde diese CVE gewählt. Sie eignet sich, um den Unterschied zwischen einer
rein versionsbasierten und einer exploit-basierten Erkennung sichtbar zu machen, da ein Dienst die
verwundbare Version tragen kann, ohne tatsächlich angreifbar zu sein.

Der zweite Dienst ist ein Apache Tomcat 9.0.30 mit CVE-2020-1938, bekannt als Ghostcat. Die Lücke
betrifft den AJP-Connector, der bei Tomcat standardmäßig auf Port 8009 lauscht. Über manipulierte
AJP-Anfragen lassen sich Dateien aus der Webanwendung lesen, etwa aus dem geschützten
WEB-INF-Verzeichnis. Unter zusätzlichen Voraussetzungen kann daraus auch eine Codeausführung werden.
Der Dienst steht im Use Case für eine Schwachstelle, die nicht über den normalen HTTP-Port, sondern
über einen separaten Protokoll-Port erreichbar ist.

Der dritte Dienst ist eine Grafana-Instanz in der Version 8.3.0 mit CVE-2021-43798. Auch hier handelt
es sich um einen Path-Traversal-Fehler. Über die Plugin-Pfade lassen sich beliebige lokale Dateien
auslesen, beispielsweise Konfigurationsdateien oder Systemdateien. Die Lücke ist ohne
Authentisierung ausnutzbar und eignet sich daher gut als klar nachweisbarer, exploit-basierter
Treffer.

Der vierte Dienst ist ein Jenkins-Server in der Version 2.138.1 mit CVE-2018-1000861. Diese Lücke
betrifft das zugrundeliegende Stapler-Web-Framework und erlaubt unter Umständen die Ausführung von
Code bereits ohne vorherige Anmeldung. Sie steht für eine besonders kritische Schwachstelle, die in
der Praxis als hochriskant gilt und sich zuverlässig über die Version erkennen lässt.

\subsubsection{Fehlkonfigurationen auf Netzwerkebene}

Die zweite Klasse besteht aus typischen Fehlkonfigurationen, wie sie in realen Umgebungen häufig
auftreten. Sie sind ebenfalls über das Netzwerk sichtbar, beruhen aber nicht auf einer
Software-CVE, sondern auf einer unsicheren Einstellung.

B1 ist ein FTP-Dienst auf Basis von vsftpd, der mit einem anonymen Zugang und mit schwachen
Zugangsdaten vorgesehen ist. Schwache Standard-Zugangsdaten wie \texttt{admin:admin} oder
\texttt{guest:guest} sollen einen einfachen, unautorisierten Zugriff ermöglichen. Der Dienst steht
für die klassische Kombination aus offenem Dienst und schwachen Anmeldedaten.

B2 ist eine absichtlich geschwächte SSH-Konfiguration. Vorgesehen sind veraltete und schwache
kryptografische Verfahren, also unsichere Ciphers, MAC- und Key-Exchange-Algorithmen, sowie
freizügige Anmeldeoptionen. Damit lässt sich prüfen, ob die Werkzeuge schwache kryptografische
Einstellungen eines ansonsten aktuellen SSH-Dienstes erkennen.

B3 sind anonyme Samba-Freigaben. Über das SMB-Protokoll werden Freigaben bereitgestellt, auf die ohne
Authentisierung zugegriffen werden kann. Dadurch sind Dateien für jeden im Netzwerk lesbar, was ein
typisches Datenleck darstellt.

B4 ist ein NFS-Export mit der Option \texttt{no\_root\_squash}. Diese Option bewirkt, dass ein
entfernter Benutzer mit Root-Rechten auch auf dem Server als Root agiert. In Kombination mit einem
offenen Export ist das eine schwerwiegende Fehlkonfiguration, die einen weitreichenden Zugriff auf
das Dateisystem ermöglichen kann.

B5 ist eine Apache-Instanz auf dem Standard-Port mit mehreren Fehlkonfigurationen. Vorgesehen sind
eine offene Status-Seite sowie eine erreichbare phpinfo-Ausgabe. Solche Seiten geben interne
Informationen über das System und die Konfiguration preis und erleichtern einem Angreifer die
Vorbereitung weiterer Schritte.

B6 ist ein rlogin-Dienst auf Basis des rsh-Servers. Dieser Dienst überträgt Anmeldedaten im Klartext
und gilt seit Langem als unsicher. Er ist als Beispiel für einen veralteten Dienst gedacht, der über
ein eigenes Protokoll erreichbar ist und sich nur schwer durch moderne, HTTP-orientierte Werkzeuge
erfassen lässt.

B7 ist ein TFTP-Dienst, der über UDP auf Port 69 erreichbar ist. TFTP erlaubt einen einfachen
Dateitransfer ohne Authentisierung. Der Dienst wurde bewusst über UDP vorgesehen, weil viele
Netzwerk-Scanner auf TCP fokussiert sind und UDP-Dienste leicht übersehen. Damit eignet er sich, um
eine bekannte Grenze netzwerkbasierter Scans zu prüfen.

\subsubsection{Host-lokale Schwachstellen}

Die dritte Klasse besteht aus Schwächen, die sich nicht über das Netzwerk zeigen, sondern erst nach
einer Anmeldung auf dem System sichtbar werden. Sie richten sich an die credentialed Perspektive (A2)
und an die Host-Audit-Perspektive (A3).

C1 betrifft sensible Dateien, die für alle Benutzer schreibbar sind (world-writable). Vorgesehen sind
eine schreibbare Sicherungskopie von \texttt{/etc/shadow} (\texttt{/etc/shadow.backup}), ein
Backup-Archiv unter \texttt{/root} und eine Log-Datei mit vertraulichem Inhalt. Besonders kritisch ist
die schreibbare Shadow-Kopie, da sich darüber Passwort-Hashes abgreifen oder manipulieren lassen.

C2 sind gesetzte SUID-Bits auf Standardwerkzeugen wie \texttt{cp}, \texttt{find} und \texttt{vim}.
Das SUID-Bit führt dazu, dass ein Programm mit den Rechten seines Eigentümers läuft, in der Regel
also als Root. Bei Standardwerkzeugen lässt sich das über bekannte Techniken zur
Rechteausweitung ausnutzen, sodass ein normaler Benutzer Root-Rechte erlangen kann.

C3 sind unsichere Einträge in der sudoers-Konfiguration. Vorgesehen sind Regeln mit der
NOPASSWD-Option, die bestimmten Konten erlauben, Befehle ohne Passwort als Root auszuführen. In
Kombination mit den genannten Werkzeugen führt das zu einer trivialen Rechteausweitung.

C4 sind schwache lokale Benutzerkonten. Vorgesehen sind mehrere Konten mit einfachen, erratbaren
Passwörtern und einer interaktiven Shell. Sie stehen für die häufige Schwäche, dass ein System zwar
aktuell gepatcht ist, aber über schwache Anmeldedaten verfügt.

\section{Evaluationsmethodik}
\label{sec:uc1-methodik}

Um die Werkzeuge objektiv zu vergleichen, braucht es zwei Dinge: einen verlässlichen Maßstab dafür,
welche Schwachstellen tatsächlich vorhanden sind, und klar definierte Kennzahlen für die Bewertung
der Funde.

\subsection{Verifizierte Referenzbasis (Ground Truth)}
\label{subsec:uc1-groundtruth}

Die in den Ansible-Rollen vorgesehenen Schwachstellen geben nur die Absicht wieder. Sie entsprechen
nicht zwangsläufig dem tatsächlichen Laufzeitzustand. Eine vorgesehene Lücke kann inaktiv sein, etwa
weil ein Konfigurations-Drop-in sie überschreibt oder ein Dienst die unsichere Einstellung nicht
übernimmt. Ein Vergleich der Werkzeuge gegen die bloße Absicht wäre daher verzerrt.

Aus diesem Grund wurde eine scanner-unabhängige Referenzbasis erstellt, in der Literatur als Ground
Truth bezeichnet. Jede geplante Schwachstelle wurde per gezieltem Direkt-Check auf dem Zielsystem
überprüft und entweder als aktiv bestätigt oder als Drift verworfen. Als Drift wird dabei eine
Abweichung zwischen dem vorgesehenen und dem tatsächlichen Zustand bezeichnet, also eine geplante
Schwachstelle, die zur Laufzeit gar nicht aktiv ist. Verwendet wurden für die Prüfung
Standardwerkzeuge wie \texttt{curl}, \texttt{find}, \texttt{grep}, \texttt{showmount},
\texttt{smbclient} und \texttt{sshd -T}. Diese Prüfungen sind gezielt und nicht-destruktiv. Sie
stellen fest, ob eine Schwachstelle real vorhanden ist, ohne das System anzugreifen.

Die Verifikation ist notwendig, weil sonst systematische Fehler entstehen. Übersieht ein Scanner eine
vorgesehene, aber gar nicht aktive Lücke, würde das fälschlich als Fehler des Scanners gewertet.
Solche Phantom-Fehler verzerren jeden Recall-Wert nach unten. Drift wird deshalb explizit
ausgewiesen und aus der Wertung ausgeschlossen.

Bei den CVE-behafteten Diensten wird zusätzlich zwischen „Version vorhanden" und „tatsächlich
ausnutzbar" unterschieden. Ein Dienst kann die verwundbare Version tragen, ohne in seiner konkreten
Konfiguration angreifbar zu sein. Wo möglich, wurde die Ausnutzbarkeit über einen gezielten,
nicht-destruktiven Check bestätigt. Bei Lücken, deren Ausnutzung das System verändern würde, bleibt
es bei der Versionsfeststellung. Diese Unterscheidung ist zentral für den Vergleich versions- und
exploit-basierter Scanner.

Die so entstandene Referenzbasis ist der alleinige Maßstab der Bewertung. Recall und Precision der
einzelnen Werkzeuge werden ausschließlich gegen diese verifizierte Liste gemessen.

\subsection{Bewertungsmetriken}
\label{subsec:uc1-metriken}

Für die Bewertung werden die beiden in Abschnitt~\ref{subsec:scan-metriken} eingeführten Kennzahlen
verwendet. Recall gibt an, welcher Anteil der real vorhandenen Schwachstellen erkannt wurde.
Precision gibt an, welcher Anteil der gemeldeten Funde tatsächlich zutrifft. Ein False Negative ist
eine übersehene, real vorhandene Schwachstelle, ein False Positive ein gemeldeter, aber nicht
vorhandener Fund.

Die Berechnung von Recall ist im Fall von UC1 nicht eindeutig, weil zwei legitime Nenner möglich
sind. Betrachtet man das vorhandene Inventar, zählt jede CVE-behaftete Software als Schwachstelle,
unabhängig davon, ob sie konkret ausnutzbar ist. Diese Sicht begünstigt versionsbasierte Scanner wie
OpenVAS und Nessus. Betrachtet man dagegen nur die tatsächlich ausnutzbaren Lücken, fallen reine
Versionstreffer heraus. Diese Sicht begünstigt exploit-basierte Werkzeuge wie Nuclei. Die Wahl des
Nenners ist damit selbst ein Ergebnis und wird in der Auswertung offengelegt.

Der reine Findings-Zähler ist als Maßstab nicht ausreichend. Ein hoher Wert kann auch durch viele
niedrig eingestufte oder nicht ausnutzbare Funde entstehen. Zur Einordnung des Schweregrads dient
daher das CVSS. Außerdem werden die Kennzahlen jeweils innerhalb einer Untersuchungsperspektive
betrachtet, da nur dort ein fairer Vergleich möglich ist.


\section{Implementierung}
\label{sec:uc1-implementierung}

Dieser Abschnitt beschreibt den konkreten Aufbau von Use Case 1 in CyberRangeCZ. Er umfasst die
Sandbox-Topologie, die automatisierte Provisionierung des Zielsystems und des Scanner-Hosts sowie
die betrieblichen Herausforderungen, die beim Aufbau gelöst wurden.

\subsection{Sandbox-Topologie}
\label{subsec:uc1-topologie}

Die Sandbox von Use Case 1 ist bewusst einfach gehalten. Sie besteht aus drei Systemen in einem
einzigen Netzwerksegment. Die folgende Topologie-Definition beschreibt die beteiligten Hosts, den
Router und das Netzwerk.

\begin{verbatim}
name: uc1-vulnscan
hosts:
  - name: target
    base_box:
      image: ubuntu-noble-x86_64
      man_user: ubuntu
    flavor: standard.small
  - name: scanner
    base_box:
      image: ubuntu-noble-x86_64
      man_user: ubuntu
    flavor: m1.medium
routers:
  - name: router
    base_box:
      image: debian-12-x86_64
      man_user: debian
    flavor: standard.small
networks:
  - name: lab-switch
    cidr: 192.168.20.0/24
net_mappings:
  - host: target
    network: lab-switch
    ip: 192.168.20.5
  - host: scanner
    network: lab-switch
    ip: 192.168.20.10
\end{verbatim}

Das Netzwerk \texttt{lab-switch} nutzt den Bereich 192.168.20.0/24. Das Zielsystem liegt auf
192.168.20.5, der Scanner auf 192.168.20.10 und der Router auf 192.168.20.1. Bewusst wird nur ein
einziges Netzwerksegment verwendet. Eine frühere Variante mit mehreren Netzen führte zu einem
Konflikt bei der Default-Route, der durch das einzelne Segment vermieden wird.

Die Flavors legen die Ressourcen der Instanzen fest. Das Zielsystem nutzt \texttt{standard.small},
der Scanner das größere \texttt{m1.medium} mit 40~GB Disk. Der Grund ist der ressourcenintensive
OpenVAS-Stack, der auf kleineren Flavors nicht lauffähig war. Auch dieser Punkt wird in
Abschnitt~\ref{subsec:uc1-betrieb} aufgegriffen.

Jede Sandbox wird in CyberRangeCZ als Satz von OpenStack-Instanzen bereitgestellt. Die
Provisionierung der Inhalte erfolgt über Ansible. Das zugehörige Playbook ruft für das Zielsystem
die Rolle \texttt{vulnerable-target} und für den Scanner die Rolle \texttt{scanner-host} auf.

\subsection{Aufbau des verwundbaren Zielsystems}
\label{subsec:uc1-zielsystem-aufbau}

Das Zielsystem wird vollständig über die Ansible-Rolle \texttt{vulnerable-target} provisioniert.
Damit ist der gesamte Aufbau reproduzierbar und versionierbar. Die vollständige Rolle sowie die
übrigen Definitionsdateien sind in Anhang~\ref{chap:quelltext-uc1} abgedruckt. Die Rolle gliedert
sich grob in drei
Teile, die den im Konzept beschriebenen Schwachstellenklassen entsprechen: die Installation und
Fehlkonfiguration netzwerk-erreichbarer Dienste, das Einrichten host-lokaler Schwächen und das
Starten der CVE-behafteten Dienste als Container.

Die Fehlkonfigurationen werden direkt über die jeweiligen Konfigurationsdateien gesetzt. Die
schwache SSH-Konfiguration wird beispielsweise als zusätzlicher Block in die \texttt{sshd\_config}
geschrieben.

\begin{verbatim}
- name: Weak SSH configuration
  blockinfile:
    path: /etc/ssh/sshd_config
    validate: '/usr/sbin/sshd -t -f %s'
    block: |
      PermitRootLogin yes
      PasswordAuthentication yes
      PermitEmptyPasswords yes
      Ciphers +aes128-cbc,3des-cbc
      MACs +hmac-sha1,hmac-md5
      KexAlgorithms +diffie-hellman-group1-sha1,diffie-hellman-group14-sha1
\end{verbatim}

Die host-lokalen Schwächen werden über einfache Dateioperationen erzeugt, etwa die world-writable
Dateien und die SUID-Bits auf Standardwerkzeugen.

\begin{verbatim}
- name: Create world-writable sensitive files
  file:
    path: "{{ item }}"
    state: touch
    mode: '0777'
  loop:
    - /etc/shadow.backup
    - /var/log/secret.log
    - /root/backup.tar.gz

- name: Set dangerous SUID bits
  file:
    path: "{{ item }}"
    mode: 'u+s'
  loop:
    - /bin/cp
    - /usr/bin/find
    - /usr/bin/vim.basic
\end{verbatim}

Die unsicheren sudoers-Regeln werden als eigene Datei unter \texttt{/etc/sudoers.d} abgelegt und beim
Schreiben mit \texttt{visudo} validiert.

\begin{verbatim}
- name: Insecure sudoers configuration
  copy:
    dest: /etc/sudoers.d/insecure
    content: |
      test ALL=(ALL) NOPASSWD: /bin/bash
      admin ALL=(ALL) NOPASSWD: ALL
      backup ALL=(ALL) NOPASSWD: /usr/bin/vim, /usr/bin/find
    mode: '0440'
    validate: 'visudo -cf %s'
\end{verbatim}

Die vier CVE-behafteten Dienste werden als Docker-Container bereitgestellt, da die verwundbaren
Versionen nicht als Standard-Ubuntu-Pakete verfügbar sind. Die Rolle installiert dafür Docker und
startet die Container in einer Schleife.

\begin{verbatim}
- name: Verwundbare CVE-Container starten
  community.docker.docker_container:
    name: "{{ item.name }}"
    image: "{{ item.image }}"
    ports: "{{ item.ports }}"
    restart_policy: unless-stopped
    state: started
    pull: true
  loop:
    - name: vuln-apache
      image: "httpd:2.4.49"
      ports: ["8080:80"]
    - name: vuln-tomcat
      image: "tomcat:9.0.30"
      ports: ["8090:8080", "8009:8009"]
    - name: vuln-grafana
      image: "grafana/grafana:8.3.0"
      ports: ["3000:3000"]
    - name: vuln-jenkins
      image: "jenkins/jenkins:2.138.1"
      ports: ["8085:8080"]
\end{verbatim}

Da Tomcat und Jenkins als JVM-Container vergleichsweise viel Arbeitsspeicher benötigen, richtet die
Rolle zusätzlich eine Swap-Datei ein.

Nicht alle vorgesehenen Schwachstellen werden im Betrieb tatsächlich aktiv. Die in der Rolle
vorgesehene verwundbare PHP-Anwendung samt MySQL-Datenbank wird nicht eingerichtet, da der zuständige
MySQL-Task eine unter MySQL 8.0 nicht mehr gültige Syntax verwendet. Die Anwendung läuft dadurch auf
Port 80 ins Leere. Die Folgen für die Bewertung werden in Abschnitt~\ref{sec:uc1-ergebnisse} als
Drift behandelt.


\subsection{Scanner-Deployment}
\label{subsec:uc1-scanner-aufbau}

Der Scanner-Host wird über die Rolle \texttt{scanner-host} provisioniert. Er stellt die
Netzwerk-Scanner sowie unterstützende Werkzeuge bereit. Zunächst installiert die Rolle
Basiswerkzeuge wie \texttt{nmap} und anschließend die Docker-Engine aus dem offiziellen
Docker-Repository.

OpenVAS wird nicht über Pakete, sondern als Greenbone-Community-Container-Stack betrieben. Die Rolle
lädt dazu die \texttt{compose.yaml} von Greenbone und startet den Verbund. Eine wichtige
Designentscheidung betrifft den Start dieses Stacks. Der erste Lauf lädt umfangreiche Images und
Feeds (CVE, SCAP, VTs), was viele Minuten dauert. Würde die Rolle darauf warten, liefe die
Sandbox-Allokation in einen Timeout. Der Pull und der Start laufen daher als Hintergrundjob ohne
Warten.

\begin{verbatim}
- name: Pull and start Greenbone (non-blocking background job)
  shell: |
    docker compose -f /opt/greenbone/compose.yaml pull
    docker compose -f /opt/greenbone/compose.yaml up -d
  args:
    chdir: /opt/greenbone
  async: 3600
  poll: 0
\end{verbatim}

Nuclei wird als Binary in einer festen Version installiert, \texttt{nmap} steht über die
Basiswerkzeuge zur Verfügung. Der dritte Scanner, Nessus Essentials, wird nicht automatisiert
ausgerollt, sondern manuell installiert, da er eine Registrierung und Lizenzaktivierung erfordert.
Diese Lizenzprüfung verhinderte später jeden verwertbaren Scan und wird in der Diskussion
(Abschnitt~\ref{sec:uc1-diskussion}) aufgegriffen. Der interaktive Zugang
zum Scanner wird über die Rolle \texttt{user-access} mit einem definierten Konto und sudo-Rechten
eingerichtet.

\subsection{Betriebliche Herausforderungen und Lösungen}
\label{subsec:uc1-betrieb}

Der Aufbau der Umgebung brachte einige nicht-triviale Probleme mit sich, die vor allem die
Automatisierung (K2), die Wartbarkeit (K7) und die Reproduzierbarkeit (K1) betreffen.

Unter Ubuntu 24.04 sind die klassischen OpenVAS-Pakete nicht mehr verfügbar, weshalb der Scanner wie
oben beschrieben als Container-Stack betrieben wird. Die dreifach verschachtelte Virtualisierung aus
physischem Host, OpenStack-VM und Sandbox-Instanzen setzte zusätzlich mehrere Ressourcengrenzen. Der
Scanner benötigte einen größeren Flavor, OpenStack deckelte die buchbare Disk künstlich, was über
eine Anpassung der Nova-Konfiguration gelöst wurde, und der ressourcenintensive OpenVAS-Stack lief
erst mit einer zusätzlichen Swap-Datei stabil.

Für den reproduzierbaren Betrieb war außerdem das Import-Verhalten der Plattform wichtig. Die
Sandbox-Definition friert den Repository-Stand beim Import ein, sodass Änderungen am Code erst nach
einem erneuten Import wirksam werden. Nach einem Neustart des Hosts müssen die verschachtelten
Schichten zudem manuell in fester Reihenfolge wieder hochgefahren werden.

Diese Punkte sind kein Mangel der Untersuchung, sondern Erfahrungswerte, die in die Bewertung der
Plattform einfließen.


\section{Ergebnisse}
\label{sec:uc1-ergebnisse}

Die Werkzeuge wurden gegen das Zielsystem ausgeführt und ihre Funde mit der verifizierten
Referenzbasis aus Abschnitt~\ref{subsec:uc1-groundtruth} abgeglichen. Dieser Abschnitt stellt zuerst
einen Überblick der Funde vor, ordnet sie dann je Schwachstelle in einer Trefferbilanz ein und
arbeitet schließlich die systematischen blinden Flecken der Werkzeugklassen heraus.

\subsection{Überblick der Funde}
\label{subsec:uc1-funde}

Vor der eigentlichen Bewertung wurde der Ist-Zustand des Zielsystems verifiziert. Dabei zeigten sich
zwei Abweichungen vom vorgesehenen Zustand. Die verwundbare PHP-Anwendung samt Datenbank war wie in
Abschnitt~\ref{subsec:uc1-zielsystem-aufbau} beschrieben nicht aktiv. Zusätzlich waren die schwachen
SSH-Anmeldeoptionen (\texttt{PasswordAuthentication}, \texttt{PermitRootLogin},
\texttt{PermitEmptyPasswords}) zur Laufzeit unwirksam, weil ein Cloud-Init-Drop-in sie überschreibt.
Die schwachen kryptografischen Verfahren des SSH-Dienstes blieben dagegen aktiv. Diese Drift wurde
festgehalten und die betroffenen Punkte aus der Wertung herausgenommen.

Tabelle~\ref{tab:uc1-kennzahlen} fasst die Funde je Werkzeug zusammen. Die Zahlen sind nicht direkt
vergleichbar, weil die Werkzeuge unterschiedliche Perspektiven abdecken und unterschiedlich zählen.

\begin{table}[H]
\centering
\footnotesize
\renewcommand{\arraystretch}{1.25}
\begin{tabular}{p{3.0cm}p{1.3cm}p{4.0cm}p{3.2cm}}
\hline
\textbf{Werkzeug} & \textbf{Persp.} & \textbf{Methode} & \textbf{Funde} \\
\hline
Nuclei & A1 & Template-/Exploit-Matching & Recall rund 70\,\% (Exploit-Ebene) \\
OpenVAS (unauth) & A1 & aktives NVT-Probing & 41 \\
OpenVAS (credentialed) & A2 & SSH-Login und Paketabgleich & 80 (plus 38 Paket-CVEs) \\
Lynis & A3 & lokales Host-Audit & Index 63, 59 Hinweise \\
OpenSCAP/CIS & A3 & XCCDF-Compliance (CIS L1) & 218 bestanden / 125 offen \\
Trivy & A4 & statische Image-Analyse & 631 (2 von 4 Images) \\
Nessus & A1/A2 & -- & nicht lauffähig (Lizenz) \\
\hline
\end{tabular}
\caption{Funde je Werkzeug. Alle quelloffenen Werkzeuge lieferten reproduzierbare Ergebnisse, Nessus war im neu aufgebauten Setup nicht nutzbar.}
\label{tab:uc1-kennzahlen}
\end{table}

In der unauthentifizierten Perspektive (A1) lieferte OpenVAS 41 Funde über aktives Probing. Nuclei
fand auf Exploit-Ebene weniger, dafür gezielter, etwa die schwache SSH-Kryptografie und die anonyme
SMB-Freigabe. Auffällig war bei Nuclei eine hohe Quote falscher Positivmeldungen im oberen
Schweregrad: von sieben als hoch oder kritisch gemeldeten Funden waren sechs falsch.

In der credentialed Perspektive (A2) stieg OpenVAS von 41 auf 80 Funde, weil der Login den
installierten Paketstand des Hosts auslesbar macht. Dazu kamen 38 weitere paketbezogene CVEs. Dieser
Zuwachs zeigt den Wert eines authentifizierten Scans gegenüber einem reinen Außenblick.

In der Host-Perspektive (A3) vergab Lynis einen Härtungsindex von 63 und 59 Hinweise. OpenSCAP prüfte
gegen ein CIS-Profil und meldete 218 bestandene und 125 offene Prüfpunkte. In der Container-Perspektive
(A4) meldete Trivy 631 Funde, allerdings nur für zwei der vier Images, da die übrigen nicht erfasst
wurden.

\subsection{Trefferbilanz gegen die Referenzbasis}
\label{subsec:uc1-trefferbilanz}

Aussagekräftiger als die reinen Zahlen ist der Abgleich gegen die bekannten, geplanten
Schwachstellen. Tabelle~\ref{tab:uc1-trefferbilanz} zeigt für jede Schwachstelle, welches Werkzeug
sie erkannt hat.

\begin{table}[H]
\centering
\footnotesize
\renewcommand{\arraystretch}{1.25}
\begin{tabular}{p{3.2cm}p{1.1cm}p{1.6cm}p{1.6cm}p{1.0cm}p{1.3cm}p{1.0cm}}
\hline
\textbf{Schwachstelle} & \textbf{Nuclei} & \textbf{OpenVAS\newline unauth.} & \textbf{OpenVAS\newline cred.} & \textbf{Lynis} & \textbf{Open-\newline SCAP} & \textbf{Trivy} \\
\hline
Apache 2.4.49 (nicht ausnutzbar) & tw. & tw. & tw. & -- & -- & nein \\
Grafana (ausnutzbar) & ja & ja & -- & -- & -- & ja \\
Tomcat Ghostcat & nein & ja & -- & -- & -- & n.\,g. \\
Jenkins & tw. & tw. & -- & -- & -- & n.\,g. \\
FTP anonymous / weak & nein & ja & ja & -- & -- & -- \\
SSH schwache Krypto & nein & tw. & ja & ja & -- & -- \\
World-writable Dateien & nein & nein & nein & nein & ja & -- \\
SUID cp/find/vim & nein & nein & nein & nein & nein & -- \\
sudoers NOPASSWD & nein & nein & nein & nein & nein & -- \\
TFTP (UDP/69) & nein & nein & nein & ja & -- & -- \\
\hline
\end{tabular}
\caption{Trefferbilanz gegen die Referenzbasis. ja = erkannt, tw. = teilweise oder abhängig, nein = verpasst, -- = nicht im Aufgabenbereich, n.\,g. = nicht gescannt.}
\label{tab:uc1-trefferbilanz}
\end{table}

Die Matrix macht mehrere Muster sichtbar. Die ausnutzbare Grafana-Lücke wurde von mehreren Werkzeugen
zuverlässig erkannt, sowohl netzwerkbasiert als auch über die Container-Analyse. Die Apache-Lücke ist
zwar in der verwundbaren Version vorhanden, in der konkreten Konfiguration aber nicht ausnutzbar. Die
versionsbasierten Werkzeuge melden sie deshalb nur abhängig vom gewählten Maßstab. Die host-lokalen
Schwächen sind für die Netzwerk-Scanner erwartungsgemäß unsichtbar. Auffällig ist, dass die
manipulierten SUID-Bits und die unsicheren sudoers-Regeln von keinem einzigen automatisierten
Werkzeug erkannt wurden.

\subsection{Recall, Precision und blinde Flecken}
\label{subsec:uc1-blindspots}

Beim Recall wirkt sich die in Abschnitt~\ref{subsec:uc1-metriken} beschriebene Wahl des Nenners
deutlich aus. Legt man die reine Inventarsicht an, in der jede vorhandene CVE-Version zählt, schneiden
die versionsbasierten Werkzeuge gut ab, weil sie auch die nicht ausnutzbare Apache-Lücke melden.
Legt man die Sicht der tatsächlichen Ausnutzbarkeit an, fällt diese Lücke heraus, und das exploit-nahe
Vorgehen von Nuclei rückt nach vorne. Keiner der beiden Maßstäbe ist allein richtig. Erst ihre
Offenlegung macht den Vergleich fair.

Bei der Precision fällt Nuclei im oberen Schweregrad ab. Sechs von sieben als hoch oder kritisch
gemeldeten Funden waren falsch. Ein hoher Findings-Zähler ist daher ohne die Prüfung gegen die
Referenzbasis irreführend, was die in Abschnitt~\ref{subsec:uc1-metriken} getroffene Entscheidung
gegen den reinen Zähler bestätigt.

Der eigentliche Erkenntnisgewinn liegt in den blinden Flecken. Jede Werkzeugklasse übersieht
systematisch etwas Bestimmtes. Die
Netzwerk-Scanner sehen nur, was exponiert und erreichbar ist. Sie sind auf TCP fokussiert, sodass der
TFTP-Dienst über UDP durchfällt, und sie können eine verwundbare Version nicht von einer tatsächlich
ausnutzbaren Lücke trennen. Der credentialed-Scan gewinnt durch den Host-Paketstand viele zusätzliche
Funde, bleibt aber blind für das Innere der Container, deren Pakete nicht im Host vorliegen. Die
Host-Audits sehen die Systemkonfiguration von innen. OpenSCAP ist als einziges Werkzeug in der Lage,
die world-writable Dateien zu erkennen. Beide Host-Audits verpassen aber die manipulierten SUID-Bits
und die sudoers-Regel, weil ihre Prüfungen generische Richtlinien abfragen und nicht die konkret
veränderten Dateien. Trivy ist das einzige Werkzeug, das in die Container-Images hineinschaut, und
findet die Grafana-Lücke sofort. Es verpasst aber die quellcode-kompilierte Apache-Instanz.

Am deutlichsten zeigt sich die Grenze der Automatisierung bei den gezielten Host-Manipulationen. Die
SUID-Bits und die unsichere sudoers-Regel entkamen allen automatisierten Werkzeugen. Erst die manuelle
Prüfung im Rahmen der Referenzbasis deckte sie auf.


\section{Diskussion}
\label{sec:uc1-diskussion}

Die Ergebnisse lassen sich zu drei übergreifenden Befunden verdichten. Diese verbinden den Use Case
mit der Bewertung der Plattform, mit den Forschungsfragen und mit dem Begriff der Cyber-Resilienz.

\subsection{Drei übergreifende Befunde}
\label{subsec:uc1-befunde}

Der erste Befund betrifft den Unterschied zwischen einer verwundbaren Version und einer tatsächlich
ausnutzbaren Lücke. Der Apache-Dienst trägt die Version mit CVE-2021-41773, ist in seiner konkreten
Konfiguration aber nicht angreifbar und antwortet auf den Ausnutzungsversuch mit einem 403. Die
versionsbasierten Werkzeuge melden ihn dennoch als verwundbar. Erst die Laufzeitprüfung zeigt den
realen Status. Für eine Bewertung der Resilienz folgt daraus, dass reine Findings-Zahlen ohne den
Kontext der Ausnutzbarkeit irreführend sind. Eine hohe Zahl gemeldeter Schwachstellen sagt für sich
genommen wenig über das tatsächliche Risiko aus.

Der zweite Befund betrifft die Drift zwischen vorgesehenem und tatsächlichem Zustand. Die absichtlich
gesetzten schwachen SSH-Anmeldeoptionen waren zur Laufzeit unwirksam, weil ein Cloud-Init-Drop-in der
Plattform sie überschreibt. Die schwachen kryptografischen Verfahren blieben dagegen aktiv. Die
Absicht in der Infrastrukturbeschreibung entspricht also nicht automatisch dem Laufzeitzustand. Genau
deshalb ist die in Abschnitt~\ref{subsec:uc1-groundtruth} beschriebene Verifikation gegen eine
unabhängige Referenzbasis notwendig. Ohne sie wären falsche Schlüsse über die Erkennungsleistung der
Werkzeuge gezogen worden. Auch die nicht eingerichtete PHP-Anwendung gehört zu dieser Kategorie.

Der dritte Befund betrifft die Lizenzierung. Der kommerzielle Scanner Nessus lieferte kein
verwertbares Ergebnis. Jeder Scan wurde von der Lizenzprüfung blockiert, die das Host-Limit als
überschritten meldete, obwohl nur wenige Adressen gescannt worden waren und die Lizenzauslastung null
anzeigte. Auch mit der freien Universitätslizenz und mehreren Anpassungen ließ sich kein vollständiger
Scan durchführen. Die quelloffenen Werkzeuge lieferten im selben Setup vollständige und wiederholbare
Ergebnisse. Für eine Cyber Range als Forschungsumgebung, die auf wiederholte und kontrollierte
Durchläufe angewiesen ist, erwiesen sich die quelloffenen Werkzeuge damit als praktikabler.

\subsection{Komplementäre Werkzeugklassen}
\label{subsec:uc1-komplementaritaet}

Über alle Perspektiven hinweg zeigt sich, dass es keinen einzelnen besten Scanner gibt. Die Werkzeuge
bilden komplementäre Schichten. Die Netzwerk-Scanner erfassen die exponierten Dienste, der
credentialed-Scan ergänzt den Paketstand des Hosts, die Host-Audits bewerten die Systemhärtung und
die Container-Analyse schaut in die Images hinein. Erst die Kombination dieser Schichten ergibt ein
vollständiges Bild. Selbst dann bleiben gezielte Host-Manipulationen wie die SUID-Bits und die
sudoers-Regel nur über eine manuelle Prüfung auffindbar.

Für die Forschungsumgebung ist weniger die Leistung eines einzelnen Werkzeugs entscheidend als die
Tatsache, dass die Umgebung den reproduzierbaren und kontrollierten Vergleich mehrerer Werkzeugklassen
gegen eine bekannte Referenzbasis überhaupt erst möglich macht. Bezogen auf die Cyber-Resilienz aus
Abschnitt~\ref{sec:cyber-resilienz} unterstützt der Use Case vor allem die vorausschauende Dimension.
Er macht die vorhandenen Schwächen eines Systems messbar und nachvollziehbar, bevor sie in einem
realen Angriff zum Tragen kommen. Die Umgebung liefert damit die Grundlage, um solche Untersuchungen
wiederholbar durchzuführen.

\subsection{Rückbezug zu den Forschungsfragen und Kriterien}
\label{subsec:uc1-rueckbezug}

Der Use Case beantwortet mehrere der Untersuchungsfragen konkret. Er zeigt, dass sich ein
reproduzierbarer Use Case auf der gewählten Plattform modellieren lässt (UF3). Die Einbindung der
externen Werkzeuge über Ansible und Docker, etwa der OpenVAS-Container und Nuclei, belegt die
Integrationsfähigkeit der Umgebung für gängige Werkzeuge (UF4). Die Ground-Truth-Methodik zusammen
mit den dokumentierten Drift-Fällen liefert eine konkrete Bewertung der Reproduzierbarkeit (UF5). Die
Drift durch das Cloud-Init-Drop-in und der fehlgeschlagene MySQL-Task sind dabei keine Nebensächlichkeiten,
sondern zeigen, an welchen Stellen Reproduzierbarkeit in der Praxis bricht.

Auch die Bewertungskriterien aus Kapitel~\ref{chap:auswahl} werden durch den Use Case bestätigt und
geschärft. Die betrieblichen Herausforderungen aus Abschnitt~\ref{subsec:uc1-betrieb} betrafen direkt
die Reproduzierbarkeit (K1), die Automatisierung (K2) und die Wartbarkeit (K7). Konkret waren das die
Anpassung der OpenStack-Konfiguration, damit die Sandbox überhaupt allokiert werden konnte, das
Löschen und Neuimportieren der Sandbox-Definition nach jeder Änderung und der manuelle Wiederanlauf
der verschachtelten Schichten nach einem Neustart. Die Einbindung der Scanner als Container betraf die
Integrationsfähigkeit (K3). Die Plattform hat sich für den Use Case als tragfähig erwiesen, allerdings
nicht ohne diesen manuellen Aufwand. Diese Erfahrungswerte fließen in die abschließende Bewertung der
Arbeit ein.
